\chapter{Touching simplices}

Two \(d\)-simplices touch if their interiors are disjoint and their
intersection is nonempty and has affine dimension \(d-1\).
Let \(f(d)\) be the largest size of a pairwise touching family in
\(\mathbb{R}^d\).

Bagemihl conjectured that for every \(d\ge 1\), we have \(f(d)=2^d\).

\begin{theorem}
  \label{ch16theorem1}
  \lean{TouchingSimplices.zaks}
  \leanok
  For every \(d\ge 2\), there is a family of \(2^d\) pairwise touching
  \(d\)-simplices in \(\mathbb{R}^d\) together with a transversal line that hits the interior
  of every single one of them.
\end{theorem}
\begin{proof}
  \leanok
  Start with four touching triangles in the plane with a common transversal.
  For the induction step, an affine transformation puts short segments
  inside all simplices into a half-square. Reflect the configuration in
  \(x_1=x_2\), then form pyramids over the two configurations with apexes
  on opposite sides of \(\mathbb{R}^d\). This doubles the number of simplices.
  A slightly tilted antidiagonal meets the interior of each pyramid.
\end{proof}

\begin{theorem}
  \label{ch16theorem2}
  \lean{TouchingSimplices.touchingNumber_lt}
  \leanok
  For all \(d \ge 1\), we have \(f(d)< 2^{d+1}\).
\end{theorem}
\begin{proof}
  \leanok
  Let \(s\) be the number of facet hyperplanes and put \(k=d+1\).
  Completing the zero entries of each row of the sign matrix gives
  \(2^{s-k}\) sign vectors. These sets of completions are pairwise disjoint,
  and at least one of the \(2^s\) possible sign vectors is missing.
  Thus \(r2^{s-k}<2^s\), so \(r<2^k=2^{d+1}\).
\end{proof}
